% Generated from locale/tr/content/lambda-calculus/lambda-definability/primitive-recursive-functions.tex; do not hand-edit.


\olfileid[tr]{lam}{ldf}{prf}
\olsection{İlkel Özyinelemeli Fonksiyonlar \usetoken{S}{lambda definable}}

İlkel özyinelemeli fonksiyonların, temel $\Zero$, $\Succ$ ve $\Proj{n}{i}$
fonksiyonlarından bileşke ve ilkel özyinelemeyle tanımlanabilen fonksiyonlar
olduğunu hatırlayın.

\begin{lem}
  \ollabel{lem:basic}
  Temel ilkel özyinelemeli $\Zero$, $\Succ$ ve izdüşüm $\Proj{n}{i}$
  fonksiyonları !!{lambda definable}.
\end{lem}

\begin{proof}
Aşağıdaki terimler bu fonksiyonları !!{lambda define}:
\begin{align*}
  \fn{Zero} & \ident \lambd[a][\lambd[fx][x]]\\
  \fn{Succ} & \ident \lambd[a][\lambd[fx][f (a f x)]] \\
  \fn{Proj}^n_i & \ident \lambd[x_0\dots x_{n-1}][x_i].
\end{align*}
\end{proof}

\begin{lem}
  \ollabel{lem:comp} $k$ yerli $f$ fonksiyonu ile $n$ yerli
  $g_0,\dots,g_{k-1}$ fonksiyonlarının sırasıyla
  $F,G_0,\dots,G_{k-1}$ terimleri tarafından !!{lambda definable}
  olduğunu ve $h$'nin bunlardan bileşkeyle tanımlandığını varsayın. O zaman
  $h$ de !!{lambda definable}.
\end{lem}

\begin{proof}
  Şu terim $h$'yi !!{lambda define}:
  \[
  H \ident \lambd[x_0 \dots x_{n-1}][F\, (G_0 x_0 \dots
    x_{n-1}) \dots (G_{k-1} x_0 \dots x_{n-1})].
  \]
  Bunun doğrulanmasını alıştırma olarak bırakıyoruz.
\end{proof}

\begin{prob}
  $H\num{n_0}\dots\num{n_{n-1}}\red
  \num{h(n_0,\dots,n_{n-1})}$ olduğunu göstererek
  \olref[lam][ldf][prf]{lem:comp} lemasının ispatını tamamlayın.
\end{prob}

\olref{lem:comp}, $f$ ve $g_0$, \dots,~$g_{k-1}$ fonksiyonlarının ilkel
özyinelemeli olmasını gerektirmez; yalnızca tümel ve !!{lambda definable}
olmaları gerekir.

\begin{lem}\ollabel{lem:prim}
  $f$ $n$ yerli, $g$ ise $n+2$ yerli bir fonksiyon olsun; bunların sırasıyla
  $F,G$ terimleri tarafından !!{lambda definable} olduğunu ve $h$'nin $f,g$
  fonksiyonlarından ilkel özyinelemeyle tanımlandığını varsayın. O zaman $h$
  de !!{lambda definable}.
\end{lem}

\begin{proof}
  $h$'nin şöyle tanımlandığını hatırlayın:
  \begin{align*}
    h(x_1, \dots, x_n, 0) &= f(x_1, \dots, x_n)\\
    h(x_1, \dots, x_n, y+1) & = g(x_1, \dots, x_n, y,
      h(x_1, \dots, x_n, y)).
  \end{align*}
  Gayriresmî olarak ilkel özyinelemeli tanım, $g$ fonksiyonunun uygulanmasını
  $y$ kez yineler ve başlangıç değeri olarak $f(x_1,\dots,x_n)$ kullanır.
  Bu, yineleyici olarak tanımlanan Church sayılarının tanımını andırır.

  Basitlik için tek bir ek $x$ argümanı bulunan duruma ait tanımı ve ispatı
  veriyoruz. Şu terim $h$'yi !!{lambda define}:
  \begin{align*}
    H \ident & \lambd[x][\lambd[y][\fn{Snd} (y
        D \tuple{\num{0}, F x})]]
    \intertext{burada}
    D \ident & \lambd[p][\tuple{\fn{Succ} (\fn{Fst}\, p), (G x
          (\fn{Fst}\, p) (\fn{Snd}\, p))}].
  \end{align*}
  Koruduğumuz yineleme durumu bir çifttir: birinci bileşen güncel $y$, ikinci
  bileşen buna karşılık gelen $h$ değeridir. Her yineleme adımında $y$ ile
  $h$'nin yeni değerlerinden oluşan bir çift kurarız; yineleme bittiğinde
  çiftin ikinci bileşenini, yani son $h$ değerini döndürürüz. Şimdi bunun
  gerçekten ilkel özyinelemenin bir gösterimi olduğunu ispatlayalım.

  Her $n,m$ için $H\,\num n\,\num m\red\num{h(n,m)}$ olduğunu göstermek
  istiyoruz. Önce $D_n\ident\Subst{D}{\num n}{x}$ ise
  $D_n^m\tuple{\num0,F\,\num n}\red
  \tuple{\num m,\num{h(n,m)}}$ olduğunu $m$ üzerinde tümevarımla gösterelim.

  $m=0$ ise göstermek istediğimiz
  $D_n^0\tuple{\num0,F\,\num n}\red\tuple{\num0,\num{h(n,0)}}$'dır.
  Sol taraf doğrudan $\tuple{\num0,F\,\num n}$'dir. $F$, $f$'yi
  !!{lambda define}s; dolayısıyla bu çift $\tuple{\num0,\num{f(n)}}$'ye
  indirgenir. $f(n)=h(n,0)$ olduğundan sonuç
  $\tuple{\num0,\num{h(n,0)}}$'dır.

  Şimdi $D_n^m\tuple{\num0,F\,\num n}\red
  \tuple{\num m,\num{h(n,m)}}$ olduğunu varsayın. Göstermek istediğimiz
  $D_n^{m+1}\tuple{\num0,F\,\num n}\red
  \tuple{\num{m+1},\num{h(n,m+1)}}$'dır.
  \begin{align*}
    D_n^{m+1} \tuple{\num{0}, F\, \num{n}}
    & \ident D_n(D_n^{m} \tuple{\num{0}, F\, \num{n}})\\
    & \red D_n\,\tuple{\num{m}, \num{h(n, m)}}
      \text{\qquad (TV gereği)}\\
    & \ident (\lambd[p][
      \tuple{
        \fn{Succ} (\fn{Fst}\, p), (G \, \num{n} (\fn{Fst}\, p) (\fn{Snd}\, p))
    }]) \tuple{\num{m}, \num{h(n,m)}}\\
    & \redone
    \tuple{\fn{Succ} (\fn{Fst}\, \tuple{\num{m}, \num{h(n,m)}}),\\
      & \qquad (G \, \num{n} (\fn{Fst}\, \tuple{\num{m}, \num{h(n,m)}})
      (\fn{Snd}\, \tuple{\num{m}, \num{h(n,m)}}))}\\
    & \red
    \tuple{\fn{Succ}\, \num{m}, (G \, \num{n} \,\num{m} \, \num{h(n,m)})}\\
    & \red \tuple{\num{m+1}, \num{g(n, m, h(n, m))}}.
  \end{align*}
  $g(n,m,h(n,m))=h(n,m+1)$ olduğundan ispat tamamlanır.

  Son olarak şunu ele alalım:
  \begin{align*}
    H\,\num n\,\num m &\ident \lambd[x][\lambd[y][\fn{Snd} (y
        (\lambd[p][\tuple{\fn{Succ} (\fn{Fst}\, p), (G\, x\,
          (\fn{Fst}\, p)\, (\fn{Snd}\, p))}]) \tuple{\num{0}, F x})]]\\
    & \qquad\,\num n\, \num m\\
    & \red \fn{Snd} (\num m \,
    \underbrace{(\lambd[p][\tuple{\fn{Succ} (\fn{Fst}\, p), (G \,\num n\,
      (\fn{Fst}\, p) (\fn{Snd}\, p))}])}_{D_n} \tuple{\num{0}, F \num n})\\
    & \ident \fn{Snd} (\num m \, D_n\, \tuple{\num{0}, F \num n})\\
    & \red \fn{Snd} \,(D_n^m  \tuple{\num{0}, F \num n})\\
    & \red \fn{Snd} \,\tuple{\num{m}, \num{h(n, m)}} \\
    & \red \num{h(n,m)}.
  \end{align*}
\end{proof}


\begin{prop}
  Her ilkel özyinelemeli fonksiyon !!{lambda definable}.
\end{prop}

\begin{proof}
  \olref{lem:basic} gereği bütün temel fonksiyonlar !!{lambda definable};
  \olref{lem:comp} ve \olref{lem:prim} gereği de !!{lambda definable}
  fonksiyonlar bileşke ve ilkel özyineleme altında kapalıdır.
\end{proof}

